\subsection{Credential Management Costs}
\label{sec:credential-cost}

For the same permission assignments, let $N$ be the number of entities, each holding at least one service permission, $p$ the average number of service permissions per entity, and $r$ the average number of roles per entity for which separate credentials are issued. A role may cover multiple services. We count credentials held by entities, not distinct shared service-key values across the deployment.

\begin{table}[t]
\centering
\caption{Entity-held credentials under shared identity trust}
\label{tab:credential-counts}
\small
\begin{tabularx}{\columnwidth}{@{}Xcc@{}}
\toprule
Credential category & NDNSF & Separate service keys \\
\midrule
Identity-signing credentials & $N$ & $N$ \\
Service-decryption credentials & $N$ & $Np$ \\
Subtotal & $2N$ & $N+Np$ \\
\bottomrule
\end{tabularx}
\end{table}

Table~\ref{tab:credential-counts} counts identity-signing and service-decryption credentials separately before adding them. Both configurations assume one active identity-signing credential per entity; a certificate and its associated private key are treated as one managed signing credential, not as one cryptographic key. NDNSF additionally assumes one ABE authority, one set of public parameters, and one current DKEY per entity. For 100 entities with 20 permissions each, the subtotals are 200 credentials for NDNSF and 2,100 for the separate-service-key configuration. These subtotals exclude content keys, additional encryption credentials, and authority-side material; any differences in those items must also be counted in a complete comparison.

Fewer credential objects do not necessarily mean fewer bytes or network messages. Let $L_i$ be the number of attribute occurrences in entity $i$'s access policy. For a construction whose key size grows linearly with $L_i$, total DKEY size is $\Theta(\sum_i L_i)$. If each policy lists its service permissions individually, this remains $\Theta(Np)$. DKEYs exceeding a Data packet's payload capacity require segmentation~\cite{dulal2023nac-abe}; separate-key deliveries can also be batched.

Granting another permission may require issuing a replacement DKEY containing the entity's updated policy. A separate-key design may instead deliver only the additional service key. Revocation that rotates the authority's parameters requires replacement DKEYs for every entity retaining permissions under that authority. Rotating only an affected service's shared key may require updates for fewer entities.

With a separately issued credential per assigned role, the role-credential count is $Nr$, or $\Theta(Nr)$. A role certificate may also be used to sign messages, so adding another $N$ identity credentials would double-count objects unless separate identity credentials are actually maintained. DNMP and trust schemas do not require a certificate per service~\cite{nichols2019dnmp,yu2015trust}; one role may cover several services. Role-based alternatives also need encryption for descriptor confidentiality, and those decryption credentials must be included when comparing complete configurations.

Aggregation reduces the number of decryption credentials to manage. Its overall cost depends on policy size, role reuse, and permission-change frequency.

\subsection{Trust-Rule Management Costs}
\label{sec:trust-rule-cost}

Trust-rule management must be counted separately from credential objects. DNMP compiles its VerSec schema into a signed binary representation for runtime use; its three-role prototype expresses the schema in 25 lines, excluding comments~\cite{nichols2019dnmp}. Parameterized rules can cover many identities, so neither DNMP nor NDNSF necessarily needs a separate trust rule for each entity or service permission.

Let $E$ denote the number of explicitly stored role--service authorization relationships. Listing these relationships requires $\Theta(E)$ entries, but does not imply $E$ independent schema rules: patterns or shared roles may express the same assignments more compactly. In NDNSF, individually listing each entity's service permissions requires $Np$ permission entries at the Controller, or $\Theta(Np)$, even though there are only $N$ DKEY objects. With fixed message formats and naming conventions, baseline signing rules in either design need not grow with $N$. NDNSF keeps service-permission lists in ABE policies; role-based designs can also change membership through credentials or authorization state without changing the schema. Changing a role's allowed actions may instead require a schema update.

For an update that delivers a complete $B$-byte signed schema to $V$ validators, the cumulative bytes received are $\Theta(VB)$. This is not the total network-link traffic, which depends on multicast, caching, and topology. A complete cost comparison must also include schema compilation, credential or DKEY updates, and permission-state distribution. Compilation and runtime validation costs must be measured separately; they cannot be inferred from credential or rule counts.

Keeping service-permission lists outside baseline signing rules does not eliminate permission-management costs; NDNSF manages them through ABE policies and keys. The invocation experiments do not measure key-distribution, schema-management, or revocation costs.
